% !PRINTED 449
\weberpage{449}
Zwischen diesen drei Polen besteht aber ein wesentlicher Unterschied. Wenden wir nämlich die Substitution $\omega$ darauf an, so bleibt der erste von ihnen ungeändert. Er ist also mindestens sechszählig und gehört zu den Substitutionen
\[
1, \chi, \chi^2, \omega, \omega\chi, \omega\chi^2.
\]
Die beiden anderen Pole, nämlich
\weq{2}{x_1 = \varrho x_2 = \varrho^2 x_3, \quad x_1 = \varrho^2 x_2 = \varrho x_3,}
gehen durch die Substitution $\omega$ in einander über, und wir müssen nachweisen, dass sie durch jede andere Substitution der Gruppe, ausser $\chi, \chi^2$, verändert werden, dass sie also nur dreizählig sind. Bezeichnen wir für den Augenblick die beiden Punkte (2) mit $\pi$ und $\pi'$, so können wir zunächst sagen, dass alle Substitutionen aus $G_{168}$, welche $\pi$ ungeändert lassen, eine Gruppe $G_\pi$ bilden
müssen, und zwar einen Theiler von $G_{168}$; die Gruppe $G_\pi$ enthält ihrerseits die cyklische Gruppe 1, $\chi$, $\chi^2$.

Ist $\sigma$ irgend eine Substitution von $G_\pi$, so kommt also $\pi$ unter den Polen von $\sigma$ vor. Daraus folgt, dass $\sigma$ weder vom $7^{\mathrm{ten}}$ noch vom $4^{\mathrm{ten}}$ Grade sein kann. Denn die Coordinaten der Pole von diesen zwei Graden sind rationale Functionen von $\varepsilon, i$, während die Coordinaten von $\pi$ die davon unabhängige Irrationalität $\varrho$ enthalten (Bd. I, \S.\,134, vergl. den Nachtrag am Ende dieses Bandes). Dass aber $\sigma$ auch nicht vom $2^{\mathrm{ten}}$ Grade sein kann, ergiebt sich daraus, dass $\pi$ auf keiner der Hauptaxen liegt. Denn läge es auf einer Hauptaxe, so müsste, weil $\varrho$ nicht rational durch $\varepsilon$ ausgedrückt werden kann, wie aus (2) mittelst der Gleichung
\[
1 + \varrho + \varrho^2 = 0
\]
abzuleiten ist, nach \S.\,118, (7) eine der 21 Relationen bestehen
\[
\begin{aligned}
\alpha &= \beta\varepsilon^{r+3} &&= \gamma\varepsilon^{-2r+1}, \\
\alpha &= \beta\varepsilon^{-3r+3} &&= \gamma\varepsilon^{-r+1}, \\
\alpha &= \beta\varepsilon^{2r+3} &&= \gamma\varepsilon^{3r+1},
\end{aligned}
\]
von denen offenbar keine möglich ist. Die Gruppe $G_\pi$ enthält also ausser der identischen nur Substitutionen $3^{\mathrm{ten}}$ Grades, und der Grad von $G_\pi$ muss eine Potenz von 3 sein. Da aber 168 nicht durch 9 theilbar ist, so muss der Grad von $G_\pi$ gleich 3 sein. Hieraus folgt:

5. \gesperrt{Es giebt 56 dreizählige Pole, die zu je zweien zu derselben Substitution gehören, und sich danach in 28 Paare ordnen. Jeder dieser Pole}
%CONT>>
